% CAPITULO 14 - MODELO DE BASE DE DATOS (ERD)

\chapter{Modelo de Base de Datos (ERD)}\label{ch:modeloerd}

El diagrama Entidad-Relación (ERD) representa el modelo de datos relacional de FisioManager. A diferencia del diagrama de clases conceptuales (Capítulo~\ref{ch:diagramaclases}), que modela el dominio del problema, el ERD define la estructura física de persistencia: tablas, columnas, tipos de datos, claves primarias, claves foráneas y restricciones de integridad.

\section{Tablas del modelo}

El modelo consta de \textbf{20 tablas} organizadas en cinco dominios funcionales:

\subsection{Dominio de Usuarios}

\begin{itemize}
  \item \codeinline{usuarios}: tabla central que almacena los datos comunes a todos los roles del sistema (nombre, email, contraseña hasheada, rol, avatar, teléfono). El campo \codeinline{rol} es un enumerado con valores \codeinline{PACIENTE}, \codeinline{TERAPEUTA} y \codeinline{ADMIN}.

  \item \codeinline{pacientes}: extensión de \codeinline{usuarios} para el perfil de paciente. Almacena edad, condición clínica, fecha de nacimiento, género, dirección y contacto de emergencia. La clave primaria es \codeinline{usuario\_id} (FK hacia \codeinline{usuarios}), implementando una relación 1:1 de especialización.

  \item \codeinline{terapeutas}: extensión de \codeinline{usuarios} para el perfil de terapeuta. Incluye especialidad, cédula profesional (única), años de experiencia, horario de atención y máximo de pacientes asignables.

  \item \codeinline{administradores}: extensión de \codeinline{usuarios} para el perfil de administrador. Define nivel de acceso (TOTAL o PARCIAL) y área de responsabilidad.
\end{itemize}

\subsection{Dominio Clínico}

\begin{itemize}
  \item \codeinline{historial\_clinico}: expediente clínico del paciente con antecedentes personales y familiares, alergias, medicamentos actuales y estado (ACTIVO/CERRADO). Cada paciente tiene exactamente un historial.

  \item \codeinline{diagnosticos}: diagnósticos registrados por terapeutas, vinculados al historial clínico. Incluyen código CIE-10, descripción, fecha, severidad (LEVE, MODERADO, SEVERO) y estado activo.

  \item \codeinline{notas\_clinicas}: notas de evolución, interconsulta o alta que el terapeuta registra durante las consultas. Vinculadas al historial clínico con fecha y tipo.

  \item \codeinline{planes\_tratamiento}: planes terapéuticos que vinculan un paciente, un terapeuta y un diagnóstico. Definen título, objetivo, fechas y estado (ACTIVO, PAUSADO, COMPLETADO, CANCELADO).

  \item \codeinline{fases\_tratamiento}: fases secuenciales que componen un plan de tratamiento. Cada fase tiene nombre, número de orden, objetivo específico, duración estimada y estado.
\end{itemize}

\subsection{Dominio de Rutinas y Ejercicios}

\begin{itemize}
  \item \codeinline{rutinas}: plantillas de rutinas de ejercicio. El campo \codeinline{paciente\_id} es nullable: cuando es \codeinline{NULL}, la rutina es un template de biblioteca; cuando tiene valor, es una rutina asignada a un paciente específico. El tipo puede ser \codeinline{TRATAMIENTO} o \codeinline{RELAJACION}.

  \item \codeinline{ejercicios}: catálogo de ejercicios reutilizables con título, descripción, tipo (FISICO/RESPIRACION), parte del cuerpo objetivo, y URLs opcionales de video e imagen demostrativa.

  \item \codeinline{parametros\_ejercicio}: tabla intermedia que vincula ejercicios con rutinas, agregando parámetros contextuales: orden dentro de la rutina, duración en minutos, repeticiones, descanso en segundos e instrucciones adicionales.

  \item \codeinline{asignaciones\_rutina}: registro de la asignación de una rutina template a un paciente por parte de un terapeuta. Incluye fechas de inicio y fin, estado (ACTIVA, PAUSADA, COMPLETADA, CANCELADA) y notas del terapeuta.

  \item \codeinline{sesiones\_ejecucion}: registro de cada ejecución de una rutina asignada. Almacena fecha de inicio/fin, estado, porcentaje de progreso, ejercicios completados y tiempo total.
\end{itemize}

\subsection{Dominio de Seguimiento}

\begin{itemize}
  \item \codeinline{feedback}: feedback post-sesión del paciente. Incluye nivel de dolor (1--10 con restricción \codeinline{CHECK}), estado emocional (enumerado de 5 valores), URL de audio opcional y comentario de texto.

  \item \codeinline{analisis\_ia}: resultado del procesamiento de IA sobre el feedback. Relación 1:1 con \codeinline{feedback} (FK única). Almacena transcripción del audio, resumen, sentimiento detectado, alertas, recomendaciones, score de confianza (0--1) y versión del modelo utilizado.

  \item \codeinline{modulos\_bienestar}: registro diario de métricas de bienestar del paciente: horas de sueño, nivel de estrés (1--10), nivel de energía (1--10), hidratación y notas personales.
\end{itemize}

\subsection{Dominio de Agenda}

\begin{itemize}
  \item \codeinline{citas}: citas programadas entre paciente y terapeuta. Vinculadas a una sala física. Incluyen estado (PROGRAMADA, CANCELADA, COMPLETADA), notas y flag de recordatorio enviado.

  \item \codeinline{salas}: catálogo de salas o consultorios disponibles con nombre, capacidad, equipamiento, piso y estado de disponibilidad.

  \item \codeinline{recordatorios}: recordatorios asociados a citas, enviados a usuarios por distintos canales (EMAIL, SMS, PUSH). Registran mensaje, fecha programada, estado de envío, fecha real de envío e intentos realizados.
\end{itemize}

\section{Diagrama ERD}

\begin{sidewaysfigure}
  \centering
  \includegraphics[width=\textheight,height=\textwidth,keepaspectratio]{FisioManager_ERD.png}
  \caption{Modelo de Base de Datos ERD de FisioManager: 20 tablas con claves primarias, claves foráneas y restricciones.}
  \label{fig:modelo-erd}
\end{sidewaysfigure}

\section{Convenciones aplicadas}

\begin{itemize}
  \item \textbf{Claves primarias (PK):} todas las tablas usan \codeinline{id} autoincremental como clave primaria, excepto las tablas de especialización (\codeinline{pacientes}, \codeinline{terapeutas}, \codeinline{administradores}) que usan \codeinline{usuario\_id} como PK y FK simultáneamente.

  \item \textbf{Claves foráneas (FK):} se definen restricciones de integridad referencial en todas las relaciones. Las FK siguen la convención \codeinline{tabla\_id} (e.g., \codeinline{paciente\_id}, \codeinline{terapeuta\_id}).

  \item \textbf{Timestamps:} las tablas principales incluyen \codeinline{created\_at} y \codeinline{updated\_at} para auditoría temporal.

  \item \textbf{Enumerados:} los campos con valores acotados se definen como \codeinline{ENUM} para garantizar integridad a nivel de base de datos.

  \item \textbf{Restricciones CHECK:} campos numéricos acotados (nivel de dolor 1--10, score de confianza 0--1) incluyen restricciones \codeinline{CHECK} para validación a nivel de base de datos.

  \item \textbf{Unicidad:} campos como \codeinline{email} en \codeinline{usuarios} y \codeinline{cedula\_profesional} en \codeinline{terapeutas} tienen restricción \codeinline{UNIQUE}. La relación 1:1 entre \codeinline{feedback} y \codeinline{analisis\_ia} se garantiza con FK única.
\end{itemize}

\section{Código PlantUML del diagrama}

El código fuente se entrega como \codeinline{12\_modelo\_erd.puml} en el directorio \codeinline{puml/}. Se regenera con:

\begin{lstlisting}[language=bash]
plantuml -tpng 12_modelo_erd.puml
\end{lstlisting}

El código fuente completo se incluye en el Anexo~\ref{ch:anexos}.
